\begingroup\fontsize{9.2}{10.0}\selectfont\setlength{\parskip}{0.10em}\setlist[itemize]{itemsep=0.10em,topsep=0.10em}
\subsection*{Claim 1 continued. Calibration also helps on in-distribution domains in \SP{}.}\reportanchor{sec:in-distribution-domain-calibration}{Calibration on in-distribution domains}
\textbf{Main findings.}
\begin{itemize}
\item \textbf{Jaccard calibration improves confidence accuracy on the in-distribution domains privacy, fraud and defamation in \SP{} with thinking on.}
Jaccard \armB{} compared with Jaccard \armA{} improves ECE in each model ($p=.020$; Table~\ref{tab:in-distribution-domain-calibration}).
\item \textbf{The same in-distribution domains in \DB{} show uncertain improvements.}
Both mean ECE changes favor Jaccard \armB{} over Jaccard \armA{}. Adjusted p-values are .542 and .899; both 95\% intervals include zero.
\end{itemize}
\reporttable{tab:in-distribution-domain-calibration}{Jaccard calibration on in-distribution domains, with thinking on.} Arrows show \armA{} followed by \armB{}. Domains receive equal weight. Seeds are shared within each model: nine for Qwen3.5-9B and three for GPT-OSS-120B. Bold entries pass the adjusted test. The rows cover 75 conversations in \SP{} and 220 in \DB{}.
\par\smallskip{\fontsize{8.5}{9.7}\selectfont\setlength{\tabcolsep}{3pt}
\begin{tabularx}{\linewidth}{@{}X>{\columncolor{blue!4}}r>{\columncolor{blue!4}}r>{\columncolor{blue!4}}r>{\columncolor{blue!4}}r>{\columncolor{orange!6}}r>{\columncolor{orange!6}}r>{\columncolor{orange!6}}r>{\columncolor{orange!6}}r@{}}\toprule
 & \multicolumn{2}{c}{\shortstack{Qwen3.5-9B\\SafePyramid\\(in-distribution\\dataset)}} & \multicolumn{2}{c}{\shortstack{Qwen3.5-9B\\DynaBench\\(out-of-distribution\\dataset)}} & \multicolumn{2}{c}{\shortstack{GPT-OSS-120B\\SafePyramid\\(in-distribution\\dataset)}} & \multicolumn{2}{c}{\shortstack{GPT-OSS-120B\\DynaBench\\(out-of-distribution\\dataset)}}\\
Metric & \armA{} $\to$ \armB{} & $p$ & \armA{} $\to$ \armB{} & $p$ & \armA{} $\to$ \armB{} & $p$ & \armA{} $\to$ \armB{} & $p$\\\midrule
\rowcolor{teal!9}\multicolumn{9}{@{}l}{\textbf{Correctness}}\\
Rule F1 & 0.764$\,\to\,$0.773 & 1.000 & 0.805$\,\to\,$0.805 & 1.000 & 0.722$\,\to\,$0.750 & 0.123 & 0.789$\,\to\,$0.785 & 1.000\\
Precision & 0.776$\,\to\,$0.780 & \textemdash & 0.815$\,\to\,$0.815 & \textemdash & 0.767$\,\to\,$0.780 & \textemdash & 0.799$\,\to\,$0.797 & \textemdash\\
Recall & 0.767$\,\to\,$0.780 & \textemdash & 0.801$\,\to\,$0.800 & \textemdash & 0.708$\,\to\,$0.751 & \textemdash & 0.785$\,\to\,$0.780 & \textemdash\\
Jaccard similarity & 0.694$\,\to\,$0.701 & \textemdash & 0.799$\,\to\,$0.799 & \textemdash & 0.629$\,\to\,$0.656 & \textemdash & 0.783$\,\to\,$0.778 & \textemdash\\
Exact rule-set accuracy & 0.331$\,\to\,$0.335 & \textemdash & 0.782$\,\to\,$0.781 & \textemdash & 0.257$\,\to\,$0.272 & \textemdash & 0.765$\,\to\,$0.760 & \textemdash\\
Parseable answers & 0.906$\,\to\,$0.922 & \textemdash & 0.967$\,\to\,$0.963 & \textemdash & 0.974$\,\to\,$0.991 & \textemdash & 0.960$\,\to\,$0.957 & \textemdash\\
\rowcolor{teal!9}\multicolumn{9}{@{}l}{\textbf{Confidence scores and calibration}}\\
Mean confidence & 0.875$\,\to\,$0.815 & \textemdash & 0.904$\,\to\,$0.851 & \textemdash & 0.906$\,\to\,$0.846 & \textemdash & 0.936$\,\to\,$0.908 & \textemdash\\
Confidence standard deviation & 0.067$\,\to\,$0.083 & \textemdash & 0.086$\,\to\,$0.113 & \textemdash & 0.056$\,\to\,$0.066 & \textemdash & 0.045$\,\to\,$0.062 & \textemdash\\
ECE $\downarrow$ (lower is better) & \textbf{0.117$\,\to\,$0.077} & 0.020 & 0.109$\,\to\,$0.096 & 0.542 & \textbf{0.261$\,\to\,$0.187} & 0.020 & 0.136$\,\to\,$0.121 & 0.899\\
\rowcolor{teal!9}\multicolumn{9}{@{}l}{\textbf{Confidence usefulness}}\\
Error detection: AUROC & 0.754$\,\to\,$0.767 & 1.000 & 0.730$\,\to\,$0.736 & \textemdash & 0.700$\,\to\,$0.751 & 0.720 & 0.675$\,\to\,$0.644 & \textemdash\\
Error detection: average precision & 0.801$\,\to\,$0.813 & \textemdash & 0.442$\,\to\,$0.420 & \textemdash & 0.857$\,\to\,$0.877 & \textemdash & 0.398$\,\to\,$0.378 & \textemdash\\
F1 loss under selection: AURC & 0.097$\,\to\,$0.093 & 1.000 & 0.103$\,\to\,$0.098 & 1.000 & 0.195$\,\to\,$0.163 & 0.314 & 0.129$\,\to\,$0.151 & 0.778\\
F1 loss found by 20\% review & 0.309$\,\to\,$0.335 & 0.882 & 0.482$\,\to\,$0.460 & \textemdash & 0.279$\,\to\,$0.302 & 1.000 & 0.401$\,\to\,$0.392 & \textemdash\\
Confidence--F1 rank correlation & 0.441$\,\to\,$0.485 & 1.000 & 0.295$\,\to\,$0.295 & \textemdash & 0.340$\,\to\,$0.438 & 0.583 & 0.189$\,\to\,$0.161 & \textemdash\\
\rowcolor{teal!9}\multicolumn{9}{@{}l}{\textbf{Evidence support}}\\
Supported rules (percentage points) & $\geq +6.3$ & \textemdash & +1.2 & 0.542 & $\geq +0.3$ & \textemdash & +0.6 & 1.000\\
\bottomrule\end{tabularx}}
\par\smallskip

\subsection*{Calibration ablation findings}\reportanchor{sec:calibration-ablations}{Calibration ablation findings}
\begin{itemize}
\item \textbf{Binary calibration improves F1 on out-of-distribution domains in \SP{} with thinking off.} For example, \armA{}-to-\armB{} F1 improves from .583 to .660 in Qwen3.5-9B and from .104 to .390 in GPT-OSS-120B ($p=.0038$ each). On out-of-distribution domains in \DB{}, Qwen3.5-9B's binary F1 worsens by 4.5 percentage points ($p=.0038$).
\end{itemize}
\subsection*{Policy-level ablation on \SP{}, with thinking on}\reportanchor{sec:policy-level-ablation}{Policy-level ablation on \SP{}, with thinking on}
\begin{itemize}
\item \textbf{Binary calibration improves the confidence-error trend as the policy changes from L0 to L2 in each model.} Binary \armB{} compared with binary \armA{} improves the L0-to-L2 ECE change ($p=.0048$ per model), starting with worse L0 ECE (Table~\ref{tab:policy-ece-trend}).
\item \textbf{The test measures policy sensitivity across in-distribution and out-of-distribution domains in \SP{}.} Each domain in Table~\ref{tab:evaluation-panel}, plus discrimination (an out-of-distribution domain), contributes four conversations. L0/L1/L2 policies give 84 positive rows excluded from training and validation. Each checkpoint answers once per row. Policies change the rules and correct answers.
\end{itemize}
{\fontsize{8.5}{9.7}\selectfont\setlength{\tabcolsep}{4pt}
\reporttable{tab:policy-ece-trend}{Binary calibration and the ECE trend.} Lower ECE is better. \armA{}/\armB{} share four Qwen3.5-9B seeds and three GPT-OSS-120B seeds. Improvement is A's L2-minus-L0 change minus B's change. P-values are adjusted for the policy-level comparisons.
\par\begin{tabularx}{\linewidth}{@{}Xrrrr@{}}\toprule
Model & A: L0 $\to$ L2 & B: L0 $\to$ L2 & Improvement & Adjusted $p$\\\midrule
\rowcolor{blue!4}Qwen3.5-9B & .039$\to$.356 & .089$\to$.261 & \textbf{.145} & .0048\\
\rowcolor{orange!6}GPT-OSS-120B & .137$\to$.530 & .240$\to$.118 & \textbf{.515} & .0048\\\bottomrule
\end{tabularx}}
\par
\par\smallskip\par\smallskip\reportfigure{fig:calibration-example}{Saved rule predictions and confidence for SFT, A and B.} The models answer the question in Figure~\ref{fig:saved-output}. Content moderation is an out-of-distribution domain in \DB{}. Rule 1 is correct. Cells show extracted rules and confidence with thinking on. Jaccard \armA{}/\armB{} use seed 2. GPT-OSS-120B SFT was already correct.
\par{\fontsize{8.5}{9.6}\selectfont\setlength{\tabcolsep}{4pt}\begin{tabularx}{\linewidth}{@{}Xlll@{}}\toprule
Model & SFT & \armA{}: correctness & \armB{}: calibration\\\midrule
\rowcolor{blue!4}Qwen3.5-9B & No rules; .90 & No rules; .90 & \cellcolor{teal!12}No rules; .60\\
\rowcolor{orange!6}GPT-OSS-120B & Rule 1; .87 & No rules; .96 & \cellcolor{teal!12}Rule 1; .93\\\bottomrule
\end{tabularx}}\par

\endgroup
